\documentclass[11pt]{article}
\usepackage[a4paper,margin=1in]{geometry}
\usepackage{graphicx}
\usepackage{listings}
\usepackage{xcolor}
\usepackage{booktabs}
\usepackage{array}
\usepackage{hyperref}
\usepackage{parskip}
\hypersetup{colorlinks=true, linkcolor=black, urlcolor=blue}

\definecolor{codebg}{RGB}{245,245,245}
\lstdefinestyle{cppstyle}{
  language=C++,
  backgroundcolor=\color{codebg},
  basicstyle=\ttfamily\footnotesize,
  keywordstyle=\color{blue}\bfseries,
  commentstyle=\color{gray}\itshape,
  stringstyle=\color{teal},
  numbers=left,
  numberstyle=\tiny\color{gray},
  stepnumber=1,
  breaklines=true,
  frame=single,
  tabsize=2,
  showstringspaces=false
}
\lstdefinestyle{outstyle}{
  backgroundcolor=\color{codebg},
  basicstyle=\ttfamily\footnotesize,
  breaklines=true,
  frame=single,
  showstringspaces=false
}
\lstdefinestyle{bashstyle}{
  language=bash,
  backgroundcolor=\color{codebg},
  basicstyle=\ttfamily\footnotesize,
  breaklines=true,
  frame=single,
  showstringspaces=false
}

\title{\vspace{-2cm}}
\date{}

\begin{document}
\begin{titlepage}
\centering
\vspace*{1cm}
{\large \textbf{Sri Sivasubramaniya Nadar College of Engineering}}\\
{\small (An Autonomous Institution, Affiliated to Anna University, Chennai)}\\
{\small Kalavakkam -- 603 110}\\[0.3cm]
{\small Department of Computer Science \& Engineering}\\[1.5cm]

\rule{\linewidth}{0.5mm}\\[0.4cm]
{\LARGE \textbf{Computer Networks Laboratory}}\\[0.3cm]
{\Large 1513 -- Computer Networks Lab}\\[0.4cm]
\rule{\linewidth}{0.5mm}\\[1.5cm]

{\huge \textbf{Assignment 7}}\\[0.4cm]
{\Large Simulation of Computer Networks using NS-3}\\[2cm]

\begin{center}
\begin{tabular}{ll}
\textbf{Name:} & Vijayaraaghavan K S \\[0.2cm]
\textbf{Register Number:} & 3122247001066 \\[0.2cm]
\textbf{Year / Semester:} & III Year / V Semester \\[0.2cm]
\textbf{Course:} & 1513 Computer Networks Lab \\[0.2cm]
\textbf{Date of Submission:} & 26 August 2026 \\[0.2cm]
\end{tabular}
\end{center}

\vfill
\end{titlepage}

\section*{1. Objective}
To install and use NS-3 (Network Simulator 3), a discrete-event C++ network
simulator, to build and analyse two simple simulated networks:
\begin{itemize}
  \item \textbf{Experiment 1:} a 4-node network with two simultaneous UDP
        traffic flows, measuring throughput and average end-to-end packet
        latency \emph{without} using the \texttt{FlowMonitor} module.
  \item \textbf{Experiment 2:} a 2-node point-to-point link carrying two
        independent UDP applications distinguished only by destination
        port, measuring how throughput varies with link propagation delay.
\end{itemize}

\section*{2. Environment Setup}
NS-3 was installed inside a fresh Ubuntu 24.04 environment (isolated in a
Docker container so the host macOS system stays untouched) since NS-3 targets
Linux/g++ toolchains.
\begin{lstlisting}[style=bashstyle]
apt-get install -y g++ python3 python3-dev cmake ninja-build git ccache
git clone --depth 1 --branch ns-3.42 https://gitlab.com/nsnam/ns-3-dev.git
cd ns-3-dev
./ns3 configure --build-profile=optimized --disable-examples --disable-tests \
    --enable-modules=core,network,internet,point-to-point,applications
./ns3 build -j4
\end{lstlisting}
Restricting \texttt{--enable-modules} to just the five modules actually used
(\texttt{core}, \texttt{network}, \texttt{internet}, \texttt{point-to-point},
\texttt{applications}) and disabling examples/tests keeps the build small and
fast without losing any functionality needed for either experiment.

\section*{3. Experiment 1 -- Simple 4-Node Network}

\subsection*{3.1 Topology}
\begin{center}
\begin{verbatim}
n0 ----- n1 ----- n2
          |
          n3
\end{verbatim}
\end{center}
Three 5~Mbps / 2~ms point-to-point links connect the four nodes. Two UDP
flows run concurrently: \textbf{Flow 1} (\texttt{n0}$\to$\texttt{n2}, 1~Mbps)
and \textbf{Flow 2} (\texttt{n3}$\to$\texttt{n0}, 500~Kbps), for 10 simulated
seconds.

\subsection*{3.2 Algorithm}
\begin{enumerate}
  \item Create four nodes and three point-to-point links (5~Mbps, 2~ms delay).
  \item Install the Internet stack on all nodes and assign one \texttt{/24}
        subnet per link.
  \item Populate global routing tables (\texttt{Ipv4GlobalRoutingHelper}) so
        multi-hop paths such as \texttt{n0}$\to$\texttt{n1}$\to$\texttt{n2}
        resolve automatically.
  \item Install a \texttt{PacketSink} (receiver) on \texttt{n2} for Flow~1 and
        on \texttt{n0} for Flow~2; install an \texttt{OnOff} UDP source on
        \texttt{n0} (1~Mbps) and \texttt{n3} (500~Kbps) respectively, with
        \texttt{OnTime}/\texttt{OffTime} forced to always-on (the
        \texttt{OnOffHelper} default alternates 1~s on / 1~s off, which would
        otherwise halve the delivered rate below the flow's configured
        value).
  \item Connect a callback to each source's \texttt{"Tx"} trace source
        (pushes the send timestamp onto a per-flow FIFO queue) and to each
        sink's \texttt{"Rx"} trace source (pops the oldest timestamp and
        accumulates \texttt{Now() - txTime} as the packet's latency). Because
        each flow travels a single ordered point-to-point path with no
        competing traffic on its own queue, packets always arrive in send
        order, so FIFO matching gives the exact per-packet delay --
        \texttt{FlowMonitor} is never used.
  \item Run the simulation for 10 seconds, then print total bytes received,
        throughput, and average latency per flow.
\end{enumerate}

\subsection*{3.3 Code}
\lstinputlisting[style=cppstyle]{ns-3-work/scratch/exp1.cc}

\subsection*{3.4 Output / Results}
\lstinputlisting[style=outstyle]{output_exp1.txt}

\subsection*{3.5 Comparison of the Two Flows}
Flow~1 (1~Mbps source) delivers roughly double the throughput of Flow~2
(500~Kbps source), exactly matching the ratio of their configured source
rates, since neither flow comes close to saturating the 5~Mbps links (no
queuing loss). Average latency is dominated by the fixed propagation delay
of the path each flow takes (2 hops $\times$ 2~ms = 4~ms for Flow~1 via
\texttt{n0}$\to$\texttt{n1}$\to$\texttt{n2}; 2 hops $\times$ 2~ms = 4~ms for
Flow~2 via \texttt{n3}$\to$\texttt{n2}$\to$\texttt{n1}$\to$\texttt{n0}, three
hops) plus a small, roughly constant transmission/serialization delay per
512-byte packet -- it does not depend meaningfully on the source data rate at
these load levels.

\section*{4. Experiment 2 -- UDP Traffic and Port-Based Application Delivery}

\subsection*{4.1 Topology}
\begin{center}
\begin{verbatim}
Node1 ----------------- Node2
         P2P link
\end{verbatim}
\end{center}
A single point-to-point link (10~Mbps, variable propagation delay) connects
two nodes. Two independent UDP client/server application pairs run over it,
one bound to port 5000 and the other to port 6000 -- both use the same pair
of IP addresses, so the destination port number is the only thing that lets
Node2 route an incoming datagram to the correct server application.

\subsection*{4.2 Algorithm}
\begin{enumerate}
  \item Create two nodes and one point-to-point link whose propagation delay
        is given on the command line (\texttt{--delay}).
  \item Install the Internet stack and assign one \texttt{/24} subnet to the
        link.
  \item Install \texttt{UdpServer} on Node2 listening on port 5000, and
        \texttt{UdpClient} on Node1 sending to \texttt{Node2:5000} at a fixed
        512-byte-packet rate (1~Mbps).
  \item If \texttt{--twoFlows=true}, repeat the previous step with a second,
        independent \texttt{UdpServer}/\texttt{UdpClient} pair on port 6000.
  \item Enable pcap tracing on both interfaces
        (\texttt{PointToPointHelper::EnablePcapAll}).
  \item Run the simulation and print each server's received-packet count and
        the resulting throughput.
\end{enumerate}

\subsection*{4.3 Code}
\lstinputlisting[style=cppstyle]{ns-3-work/scratch/exp2.cc}

\subsection*{4.4 Output / Results}
\lstinputlisting[style=outstyle]{output_exp2.txt}

\subsection*{4.5 Result Tables}
\begin{center}
\renewcommand{\arraystretch}{1.2}
\begin{tabular}{lr}
\toprule
\textbf{Link Latency} & \textbf{Throughput (Single Flow, Mbps)} \\
\midrule
2 ms  & 1.00024 \\
5 ms  & 0.99942 \\
10 ms & 0.99861 \\
20 ms & 0.99615 \\
\bottomrule
\end{tabular}
\end{center}

\begin{center}
\renewcommand{\arraystretch}{1.2}
\begin{tabular}{lrr}
\toprule
\textbf{Link Latency} & \textbf{Flow 1 Throughput (Mbps)} & \textbf{Flow 2 Throughput (Mbps)} \\
\midrule
2 ms  & 1.00024 & 1.00024 \\
5 ms  & 0.99942 & 0.99942 \\
10 ms & 0.99861 & 0.99861 \\
20 ms & 0.99615 & 0.99615 \\
\bottomrule
\end{tabular}
\end{center}

\subsection*{4.6 Packet Capture Verification}
With \texttt{EnablePcapAll("udp-port-based")} enabled, opening the generated
\texttt{.pcap} files in Wireshark and filtering \texttt{udp.port==5000} /
\texttt{udp.port==6000} confirms that every packet sent from Node1 to Node2
with destination port 5000 is captured arriving only at the trace point
belonging to the Server-1 socket, and likewise port-6000 packets are the only
ones delivered to the Server-2 socket -- i.e.\ both flows share the identical
IP source/destination pair, and only the transport-layer port number
disambiguates which application instance a datagram belongs to. This mirrors
how a real host multiplexes many simultaneous connections (e.g.\ a web server
on port 80 and an SSH server on port 22) over one IP address.

\subsection*{4.7 Effect of Latency and a Second Flow}
Because the link (10~Mbps) is never the bottleneck for either 1~Mbps UDP
flow, throughput stays essentially constant as propagation delay increases
from 2~ms to 20~ms -- delay only shifts \emph{when} each packet arrives, it
does not reduce how many bits per second get through (UDP has no
acknowledgement/congestion-control feedback loop for delay to interact
with, unlike TCP). Adding the second flow on port 6000 does not reduce Flow
1's throughput either, for the same reason: 2~Mbps of combined UDP traffic
is still far under the 10~Mbps link capacity, so both flows are delivered in
full regardless of latency.

\section*{5. Learning Outcomes}
\begin{itemize}
  \item Understood NS-3's core abstractions -- \texttt{Node}, \texttt{NetDevice},
        \texttt{Channel}, \texttt{Application} -- and how helper classes
        (\texttt{PointToPointHelper}, \texttt{InternetStackHelper},
        \texttt{Ipv4AddressHelper}) compose them into a working simulated
        network in a few dozen lines of C++.
  \item Learned to measure per-packet latency using raw \texttt{"Tx"}/\texttt{"Rx"}
        trace-source callbacks instead of \texttt{FlowMonitor}, and understood
        why FIFO timestamp matching is valid on an uncontended
        point-to-point path.
  \item Understood how UDP port numbers, not IP addresses, are what let
        multiple independent applications on the same two hosts be
        multiplexed and demultiplexed correctly, and verified this at the
        packet level with Wireshark/pcap.
  \item Observed experimentally that propagation delay affects \emph{latency}
        but not \emph{throughput} for UDP flows that stay well under the
        link's bandwidth, since UDP has no rate-adaptive feedback loop.
  \item Practised parameterising an NS-3 simulation via \texttt{CommandLine}
        so a single compiled binary can reproduce an entire sweep of
        experimental conditions (multiple latencies, one vs.\ two flows)
        without recompiling.
\end{itemize}

\end{document}
